\subsection{HIRE}

Recent studies have shown that hallucination-related information can be identified in the hidden representations of LVLMs. Based on this observation, HIRE \citep{suo2026hallucination} proposes a representation-level hallucination manipulation framework consisting of an editor and a router.

HIRE learns to disentangle semantic and hallucination-related representations through image perturbation. Clean and perturbed images are encoded into latent representations, and a clustering objective is used to separate hallucination-related and reliable latent representations while preserving semantic information. Based on the learned representation space, HIRE derives editing directions that can either suppress or amplify hallucination behavior through a controllable editing scale.

To improve efficiency, HIRE further introduces a router that selectively applies editing to a subset of tokens. The router is trained using Dynamic Preference Optimization with randomly generated editing trajectories, enabling efficient hallucination mitigation while reducing computational cost.

Although HIRE achieves strong hallucination mitigation performance, it requires additional training for representation editing and routing. Furthermore, the relationship between perturbation-induced representations and hallucination formation remains insufficiently explained.

